% Generated deterministically from the audited Markdown source.
% Responsible author: Carptopus.
\documentclass[11pt]{article}
\usepackage[a4paper,margin=27mm]{geometry}
\usepackage[T1]{fontenc}
\usepackage[utf8]{inputenc}
\usepackage{lmodern}
\usepackage{microtype}
\usepackage{amsmath,amssymb,amsthm,mathtools}
\usepackage{booktabs}
\usepackage{needspace}
\usepackage{xcolor}
\usepackage[colorlinks=true,linkcolor=blue!55!black,citecolor=blue!55!black,urlcolor=blue!65!black]{hyperref}
\usepackage{fancyhdr}

\newtheorem{theorem}{Theorem}[section]
\newtheorem{corollary}[theorem]{Corollary}
\numberwithin{equation}{section}

\pagestyle{fancy}
\fancyhf{}
\fancyhead[L]{Carptopus}
\fancyhead[R]{Middle runner counterexamples}
\fancyfoot[C]{\thepage}
\setlength{\headheight}{14pt}
\setlength{\parindent}{0pt}
\setlength{\parskip}{0.45em}
\setlength{\emergencystretch}{4em}
\urlstyle{same}

\title{A continuous family of counterexamples\\
to the middle runner conjecture}
\author{Carptopus\\\href{mailto:carptopus@163.com}{carptopus@163.com}}
\date{27 August 2026}

\hypersetup{
  pdftitle={A continuous family of counterexamples to the middle runner conjecture},
  pdfauthor={Carptopus},
  pdfsubject={An exact continuous family of three-runner counterexamples to the middle runner conjecture as stated in the 2015 AIM problem list},
  pdfkeywords={middle runner conjecture, dynamical algebraic combinatorics, circular runners, order statistics, counterexample, piecewise integration}
}

\begin{document}
\maketitle

\begin{center}
\fcolorbox{black!35}{black!4}{\parbox{0.88\textwidth}{\centering\small
Internal candidate manuscript, version 0.1-beta. The mathematical proof and
complete manuscript have passed project-internal zero-trust audits. External
mathematical review and formal peer review are pending.}}
\end{center}

\begin{abstract}


The middle runner conjecture from the 2015 AIM workshop problem list concerns
the time averages of the ordered positions of runners moving at commensurable
nonzero speeds on a circular track. As stated, the runners may start at
arbitrary positions, and the conjecture asserts that complementary order
statistics have averages summing to one. We give a continuous family of
counterexamples with three runners. For velocities $(1,2,3)$ and initial
positions $(0,0,a)$, let $M(a)$ be the average median position over one common
period. We prove the exact formula

\nopagebreak[4]
\[
M(a)=
\begin{cases}
\displaystyle \frac12-\frac49a\left(\frac12-a\right),
  &0\leq a\leq\frac12,\\[6pt]
\displaystyle \frac12+\frac49\left(a-\frac12\right)(1-a),
  &\frac12\leq a\leq1.
\end{cases}
\]

Consequently, modulo one, the conjectured equality holds in this family only
for $a=0$ and $a=1/2$; every other offset is a counterexample. In particular,
$a=1/4$ gives ordered-position averages

\nopagebreak[4]
\[
\left(\frac{161}{576},\frac{17}{36},\frac{431}{576}\right).
\]

The proof is an elementary but complete piecewise integration. Exact rational
verification code is provided as supplementary evidence.

\end{abstract}

\section{Introduction}


The 2015 AIM workshop problem list on dynamical algebraic combinatorics asks
the following question \cite[Section~2.1]{aim2015}. Consider $n$ runners moving on a
circular track at commensurable nonzero speeds. Positions are normalized to
$[0,1)$, and the runners are explicitly not required to start at the starting
line. If

\nopagebreak[4]
\[
p_1(t)\leq p_2(t)\leq\cdots\leq p_n(t)
\]

are the ordered positions at time $t$, let $p_i$ be the average of $p_i(t)$
over one common period. Conjecture 2.1 of \cite{aim2015} states that

\nopagebreak[4]
\[
p_i+p_j=1\qquad\text{whenever }i+j=n+1.
\tag{1.1}
\]

For odd $n$, the central case says that the average position of the middle
runner is $1/2$, which gives the problem its name. The source reports that the
equal-speed case had been solved and that the general case remained open in
2015.

This note disproves (1.1) under its published hypotheses. The smallest
nontrivial odd case $n=3$ already contains a one-parameter family for which the average
middle position can be evaluated exactly. The parameter space is classified
completely: only two offsets modulo one satisfy the proposed symmetry, while
all other offsets fail it. Thus the obstruction is not an isolated numerical
configuration.

The result is deliberately limited to the fixed velocity triple $(1,2,3)$ and
the initial positions $(0,0,a)$. It is not a classification of arbitrary
three-runner systems and does not propose a corrected conjecture.

\section{Setup and main result}


For $x\in\mathbb R$, write $\{x\}=x-\lfloor x\rfloor\in[0,1)$ for its
fractional part. Fix $a\in\mathbb R/\mathbb Z$ and consider the three
positions

\nopagebreak[4]
\[
x_1(t)=\{t\},\qquad
x_2(t)=\{2t\},\qquad
x_3(t)=\{3t+a\}.
\tag{2.1}
\]

The velocities are positive, nonzero, and commensurable, and one common
period is $1$. Let

\nopagebreak[4]
\[
m_a(t)=\mathrm{med}\bigl(x_1(t),x_2(t),x_3(t)\bigr)
\]

and define

\nopagebreak[4]
\[
M(a)=\int_0^1m_a(t)\,dt.
\tag{2.2}
\]

Changing $a$ by an integer does not change the system, so we use the
representative $a\in[0,1]$.

\Needspace{0.25\textheight}
\begin{theorem}

For the system in (2.1),

\nopagebreak[4]
\[
M(a)=
\begin{cases}
\displaystyle \frac12-\frac49a\left(\frac12-a\right),
  &0\leq a\leq\frac12,\\[6pt]
\displaystyle \frac12+\frac49\left(a-\frac12\right)(1-a),
  &\frac12\leq a\leq1.
\end{cases}
\tag{2.3}
\]

Hence $M(a)=1/2$ exactly when $a\in\{0,1/2,1\}$. Equivalently, modulo one,
the middle-runner equality holds in this family exactly for
$a\in\{0,1/2\}$.

\end{theorem}

\Needspace{0.25\textheight}
\begin{corollary}

For every

\nopagebreak[4]
\[
a\in(\mathbb R/\mathbb Z)\setminus\{0,1/2\},
\]

the three-runner system (2.1) is a counterexample to Conjecture 2.1 of \cite{aim2015}.
The counterexamples form two open arcs in the parameter circle and therefore
an uncountable continuous family.

\end{corollary}

\section{Piecewise evaluation of the median}


We first assume

\nopagebreak[4]
\[
0<a<\frac12.
\tag{3.1}
\]

All changes in the relative order of the three trajectories occur at a wrap
point or at an intersection. The internal wrap points of $x_3$ are

\nopagebreak[4]
\[
A=\frac{1-a}{3},\qquad
C=\frac{2-a}{3},\qquad
F=1-\frac a3,
\]

and the internal wrap point of $x_2$ is $1/2$. The intersections of $x_1$
with $x_3$ occur at

\nopagebreak[4]
\[
B=\frac{1-a}{2},\qquad E=1-\frac a2,
\]

while the only internal intersection of $x_2$ with $x_3$ is

\nopagebreak[4]
\[
D=1-a.
\]

The first two trajectories meet only at the endpoints of the period. Under
(3.1), these events have the strict order

\nopagebreak[4]
\[
0<A<B<\frac12<C<D<E<F<1.
\tag{3.2}
\]

On each resulting open interval the relative order is fixed. Evaluating the
three affine pieces at any interior point gives the following median.

\begin{center}
\begin{tabular}{ll}
\toprule
interval & $m_a(t)$ \\
\midrule
$(0,A)$ & $2t$ \\
$(A,B)$ & $t$ \\
$(B,1/2)$ & $a+3t-1$ \\
$(1/2,C)$ & $t$ \\
$(C,D)$ & $2t-1$ \\
$(D,E)$ & $a+3t-2$ \\
$(E,F)$ & $t$ \\
$(F,1)$ & $2t-1$ \\
\bottomrule
\end{tabular}
\end{center}


Values at the finitely many intersections, wrap discontinuities, and period
endpoints do not affect the integral. Therefore

\nopagebreak[4]
\[
\begin{aligned}
M(a)={}&A^2+\frac{B^2-A^2}{2}
 +(a-1)\left(\frac12-B\right)
 +\frac32\left(\frac14-B^2\right)\\
&+\frac{C^2-1/4}{2}
 +(D^2-C^2)-(D-C)\\
&+(a-2)(E-D)+\frac32(E^2-D^2)
 +\frac{F^2-E^2}{2}+F-F^2.
\end{aligned}
\tag{3.3}
\]

Substituting $A,\ldots,F$ and collecting terms yields

\nopagebreak[4]
\[
M(a)=\frac{8a^2-4a+9}{18}
=\frac12-\frac49a\left(\frac12-a\right).
\tag{3.4}
\]

The endpoint values $a=0$ and $a=1/2$ follow either by direct subdivision or
by dominated convergence, since for all but finitely many $t$ the trajectories
and their median converge pointwise as $a$ approaches either endpoint.

For clarity, write $x_i^{(b)}$ for the trajectory in (2.1) with offset
parameter $b$; the first two trajectories do not depend on $b$. To obtain the
second half of (2.3), compare parameters $a$ and $1-a$. Away
from the finitely many wrap times, time reversal and reflection of the circle
give the multiset identity

\nopagebreak[4]
\[
\{x_1^{(a)}(1-t),x_2^{(a)}(1-t),x_3^{(a)}(1-t)\}
=
\{1-x_1^{(1-a)}(t),1-x_2^{(1-a)}(t),1-x_3^{(1-a)}(t)\}.
\]

Reflection reverses the order of three numbers but sends their median to one
minus the original median. Hence, almost everywhere,

\nopagebreak[4]
\[
m_a(1-t)=1-m_{1-a}(t).
\tag{3.5}
\]

Integration gives

\nopagebreak[4]
\[
M(a)=1-M(1-a).
\tag{3.6}
\]

Using (3.4) for $1-a\in[0,1/2]$ proves the second line of (2.3). The two
quadratic factors in (2.3) vanish only at $0,1/2,1$, completing the proof of
Theorem 2.1 and Corollary 2.2.

\section{A rational counterexample and the paired equality}


The parameter $a=1/4$ gives a compact exact witness. Integrating all three
order statistics over the same subdivision gives

\nopagebreak[4]
\[
(p_1,p_2,p_3)=
\left(\frac{161}{576},\frac{17}{36},\frac{431}{576}\right).
\tag{4.1}
\]

In particular,

\nopagebreak[4]
\[
p_2=\frac{17}{36}\neq\frac12,
\qquad
p_1+p_3=\frac{37}{36}\neq1.
\tag{4.2}
\]

The simultaneous failure in (4.2) is forced throughout the family. Each
individual positive-integer-speed trajectory is uniformly distributed over
the circle during the common period and therefore has average $1/2$.
Pointwise, sorting does not change the sum of the three positions, so

\nopagebreak[4]
\[
p_1+p_2+p_3=\frac32.
\tag{4.3}
\]

Consequently,

\nopagebreak[4]
\[
p_1+p_3=\frac32-M(a).
\tag{4.4}
\]

Thus the middle equality $p_2=1/2$ and the complementary equality
$p_1+p_3=1$ hold or fail together, exactly at the parameters classified in
Theorem 2.1.

\section{Exact verification}


The supplementary Python program
{\small\path{verification/}\allowbreak\path{verify_middle_runner_counterexample.py}} uses exact rational
arithmetic. It independently generates every wrap point and pairwise
intersection, orders the affine pieces on each interval, and integrates all
three order statistics. It checks:

\begin{itemize}
\item the witness $a=1/4$ in (4.1);
\item the common-start positive control $a=0$;
\item invariance under replacing an initial position by an integer translate; and
\item formula (2.3) for all $775$ distinct rational representatives in $[0,1]$
with denominator at most $50$. Since $0$ and $1$ represent the same point,
these are $774$ distinct parameter classes modulo one.

\end{itemize}

The program uses only the Python standard library and requires Python 3.9 or
later. From the directory containing this manuscript and the {\small\path{verification}}
subdirectory, run

\begin{verbatim}
python verification/verify_middle_runner_counterexample.py
\end{verbatim}


The finite checks are supplementary evidence and regression tests. They do
not prove the statement for all real parameters; the continuous result rests
on the piecewise argument in Section 3.

\section{Scope and literature boundary}


The contribution of this note is the exact evaluation of a complete
one-parameter family and the resulting disproof of the middle runner
conjecture under the hypotheses printed in \cite{aim2015}. The phrase "as stated" is
important only as a matter of scope: arbitrary starting positions are not an
inference made here but an explicit part of the published problem statement.

As of 27 August 2026, a targeted search through the exact problem title and wording, combinations
of the four problem-list authors with the relevant terminology, and the
arXiv, OpenAlex, and Crossref indexes did not locate a later correction or a
directly equivalent counterexample. This is not an exhaustive priority
certificate. Unindexed notes, private communications, differently named
equivalent results, and later parallel work may exist. Accordingly, this note
does not claim absolute global priority, external expert confirmation, or
peer review.

The theorem does not classify arbitrary initial phases or arbitrary velocity
triples, and it does not settle any strengthened version that imposes a
common starting position. Such variants are outside the present claim.

\section*{Acknowledgement of AI assistance}


The author used OpenAI Codex for literature triage, exact computation,
symbolic checks, proof stress testing, verification-code review, and drafting
assistance. The author reviewed the mathematical argument and assumes
responsibility for the claims and presentation.

\begin{thebibliography}{9}
\footnotesize
\raggedright

\bibitem{aim2015}
J.~Propp, T.~Roby, J.~Striker, and N.~Williams,
\emph{Problems for 2015 AIM workshop on Dynamical Algebraic Combinatorics},
29 May 2015, Section~2.1, Conjecture~2.1.
\href{https://aimath.org/pastworkshops/dac_preworkshop.pdf}{AIM problem list}.

\end{thebibliography}

\end{document}
